% SHARD-ID: RC-03C-GRADED-RAMANUJAN
% SHARD-TITLE: The Ramanujan Property for Graded Transfer Channels: Definition, Graded Quantum Ihara-Bass, Graded Harrow Expanders
% SHARD-SUMMARY: TJO asked for a rigorous definition of the Ramanujan property that starts from graded MPS, survives the continuum limit, and admits Harrow expanders as a test, suspecting that all known quantum expanders are ungraded; the suspicion is confirmed and the definition is given on the divisor of the parity-closed ring zeta.
% SHARD-SUMMARY: Graded quantum Ihara-Bass holds sector by sector, so index-two (Clifford) gradings of Harrow channels are graded Ramanujan expanders whose odd sector carries zeros on the critical circle: PGL_2(F_p) with LPS generators in the odd coset, and the six-letter Pauli channel on a qubit whose periodic non-backtracking ring norms count the points of an elliptic curve over F_5.
% SHARD-KEYWORDS: Ramanujan, Riemann hypothesis, graded bond, supertrace, divisor, quantum expander, Hashimoto, Ihara-Bass, Alon-Boppana, Clifford theory, PGL_2, LPS, Hecke, Pauli, elliptic curve
% SHARD-NOTES: none
% SHARD-SCRIPTS: scripts/graded_ramanujan.py

\section{The Ramanujan Property for Graded Transfer Channels}
\label{sec:graded-ramanujan}

\emph{Question (TJO, 2026-09-15).} Working from the MPS picture, and given that the graded
supertrace is necessary for zeta zeros, what is the correct definition of the Ramanujan property,
rigorous and meaningful in the continuum limit, with Harrow expanders as a test case? Do we
understand it in the graded case at all, or do we only have quantum expanders for ungraded bonds
and antiperiodic closures? \emph{Answer.} The suspicion is right about the literature and about
this book up to \cref{sec:zeta-conditions}: every quantum expander in sight is ungraded, and on an
ungraded bond the two closures coincide and the ring zeta has only poles
(\cref{prop:no-ungraded-zeros}). The definition is \cref{def:graded-rh-fe-ramanujan}, stated on
the divisor of the periodic ring zeta, and it is exactly stable under the continuum limit
(\cref{prop:graded-continuum-exact,prop:graded-chernoff}). Graded quantum expanders exist with a
mechanism (\cref{thm:graded-harrow}), the smallest being a qubit with Pauli letters whose graded
Ihara zeta is an elliptic curve zeta (\cref{prop:qubit-graded-zeta}). Everything in this section is
the notebook's own work of one session (author claude:fable-5.1, numerics
\texttt{scripts/graded\_ramanujan.py}, $81$ checks), with one codex gpt-6-astra prover lane
(\texttt{notes/ramanujan-graded/astra-proofs.md}: $4$ proved, $9$ corrected, $1$ refuted, a
$41$-row correction ledger, $132$ further checks) whose corrections are applied below and recorded
in the claims database; nothing here has had an independent (second-family) review.

\subsection{Positivity and the $P$-mode}

\begin{proposition}[Graded Weil--Hodge inequality]
\label{prop:graded-weil-hodge}
\claimstatus{prop:graded-weil-hodge}{sketched}
For every graded transfer channel and $n\ge1$, $|\Tr\Esec{1}^n|\le\Tr\Esec{0}^n$; hence
$\specrad(\Esec{1})\le\specrad(\Esec{0}) = \qs$, both sector spectra being conjugation closed.
Equal radii do not force a common eigenvalue: the single even letter $P$ on $\mathbb C^{1|1}$ has
even spectrum $\{1,1\}$ and odd $\{-1,-1\}$ (prover, ledger L01).
\end{proposition}
\begin{proof}
Both closures are ring norms: $\Tr\Edb^n = \sum_w|\Tr\Kr{w}|^2\ge0$ and $\str\Edb^n =
\sum_w|\Tr(P\Kr{w})|^2\ge0$ (\cref{def:graded-lattice-tensor}), i.e.\ $\Tr\Esec{0}^n\ge\pm\Tr
\Esec{1}^n$; a Perron eigenmatrix $R\ge0$ averages to the even Perron eigenmatrix $(R+PRP)/2$,
so the radius is attained evenly (positive operators need not be even, ledger L02). Proved by the
prover lane (D1). This is the finite shadow of the Hodge index inequality of
\cref{sec:ring-norm-bosonic-nogo}.
\end{proof}

\begin{proposition}[The $P$-mode]
\label{prop:p-mode}
\claimstatus{prop:p-mode}{sketched}
$\Edb(P) = P\sum_s\epsilon_s\Kr{s}\Kr{s}^\dagger$. Under parity-twisted unitality
(\cref{def:parity-twisted-tp}) $P$ is an even eigenvector with eigenvalue $\pmode$; for a graded
representation channel $\pmode = \epsilon(W_S)$, the sign character evaluated on the walk element.
The $P$-mode is structural but not trivial: for the elliptic tensor of \cref{sec:zeta-conditions}
$\pmode = 1$ (projective) or $0$ (affine), the $H^0$ pole and the functional-equation partner of
the fixed point; for the Pauli letters below $\pmode = -2$ lies inside the Hastings band and its
Ihara--Bass quadratic contributes a pole pair on the critical circle; for all letters odd
$\pmode = -\qs-1$ is the period image of the fixed point (bipartite case). For unitary letters the
Ramanujan property therefore imposes the \emph{balance condition} $|w_{\mathrm{even}}-
w_{\mathrm{odd}}|\le\lamH$ unless $\pmode\in\trivS$.
\end{proposition}
\begin{proof}
$P\Kr{s}P = \epsilon_s\Kr{s}$ gives $\Kr{s}P\Kr{s}^\dagger = \epsilon_sP\Kr{s}\Kr{s}^\dagger$; for a
representation channel $\pi(s)P\pi(s)^\dagger = \epsilon(s)P$. Proved by the prover lane (D2).
Checked on every example of \cref{num:graded-ramanujan}, including
$\sum_s\epsilon_s\Kr{s}\Kr{s}^\dagger = 1$ for the 06h tensor.
\end{proof}

\subsection{Graded quantum Ihara--Bass and graded quantum expanders}

\begin{proposition}[No zeros without a grading; no unique fixed point with only even letters]
\label{prop:no-ungraded-zeros}
\claimstatus{prop:no-ungraded-zeros}{sketched}
(i) For $P = 1$ the two closures coincide and the ring zeta $1/\det(1-\uz\Edb)$ has no zeros; a
quantum expander in the sense of \cref{def:quantum-expander} (Hastings, Harrow, the Weil--LPS
channels of \cref{sec:weil-lps}, the complexes of \cref{sec:complex-zeta}) is such a channel.
(ii) If every letter is even, the bond splits, $\str\Edb^n = \|\Psi_+-\Psi_-\|^2$, and both
$1_+$ and $1_-$ are fixed: no unique fixed point. Zeros with a unique fixed point require odd
letters, i.e.\ letters that anticommute with $P$. (iii) No nonzero positive operator is odd; odd
fixed operators can nevertheless exist when the stationary state is not unique (the identity
channel fixes every coherence), and uniqueness excludes them (ledger L29).
\end{proposition}
\begin{proof}
(i) is the definition of the ring zeta; (ii) is the second reading after
\cref{thm:cmps-twisted-supertrace} together with $\Edb(1_\pm) = 1_\pm$ for even unital letters;
(iii) $PXP = -X\ge0$ forces $X = 0$, and a Hermitian fixed operator of a finite CPTP map is a
combination of stationary states. Proved by the prover lane (D14).
\end{proof}

\begin{theorem}[Graded quantum Ihara--Bass]
\label{thm:graded-qihara-bass}
\claimstatus{thm:graded-qihara-bass}{sketched}
For the graded Hashimoto operator of \cref{def:graded-hashimoto} with $\Sig_k = \Dk\chan|_{k}$,
$n_0 = D_+^2+D_-^2$ and $n_1 = 2D_+D_-$,
\[
  \det(1-\uz\Hash_k) = (1-\uz^2)^{n_k(\Dk-2)/2}\det\bigl(1-\uz\Sig_k+(\Dk-1)\uz^2\bigr),\qquad k = 0,1,
\]
so the graded quantum Ihara zeta is $(1-\uz^2)^{-(\Dk-2)(D_+-D_-)^2/2}\prod_{\lambda\
\mathrm{odd}}(1-\lambda\uz+\qs\uz^2)\big/\prod_{\lambda\ \mathrm{even}}(1-\lambda\uz+\qs\uz^2)$,
$\qs = \Dk-1$, and for a balanced grading the $(1-\uz^2)$ factors cancel exactly.
\end{theorem}
\begin{proof}
$\Gdb\otimes1$ commutes with $\Hash$, with $\Sig\otimes1$ and with the reversal; the proof of
\cref{thm:qihara-general} in the unitary case (\cref{cor:qihara-unitary}) runs on each invariant
subspace with $n^2$ replaced by its dimension. Proved by the prover lane (D3). Verified to
$10^{-9}$ on all three examples of \cref{num:graded-ramanujan}.
\end{proof}

\begin{corollary}[Band if and only if circle, sector by sector]
\label{cor:graded-band-circle}
\claimstatus{cor:graded-band-circle}{sketched}
For a Hermitian graded channel (unitary letters closed under adjoint, equal weights) each
adjacency eigenvalue $\lambda$ of a sector gives two edge eigenvalues with $\edgeev_+\edgeev_- =
\qs$, both of modulus $\sqrt\qs$ iff $|\lambda|\le2\sqrt\qs$. Hence: every point of the
\emph{retained net} divisor of the graded quantum Ihara zeta lies on $|\uz| = \qs^{-1/2}$ iff
$|\lambda|\le2\sqrt\qs$ for every $\lambda\notin\{\pm(\qs+1)\}$ with nonzero net multiplicity
$m_0(\lambda)-m_1(\lambda)$. A bound on both entire sectors implies the divisor statement; the
converse is false without a no-cancellation hypothesis (prover, D4: a primitive degree-$14$ channel
with even $\{14,10\}$, odd $\{10,6\}$ has the mode $10$ outside the band, cancelled, and reduced
zeta $(1-6\uz+13\uz^2)/((1-\uz)(1-13\uz))$). Zeros are the net-odd quadratics and poles the net-even
ones; the trivial factors are $\{1,\qs\}$ from $\lambda = \qs+1$, $\{-1,-\qs\}$ from $\lambda =
-(\qs+1)$ when present, and $\pm1$ from the $(1-\uz^2)$ factors; zeta points are the reciprocals of
edge eigenvalues.
\end{corollary}
\begin{proof}
\cref{thm:graded-qihara-bass} and the root analysis of \cref{obs:rh-iff-ramanujan-channel}, applied
to each sector and then to the net multiplicities; $\Sig_1$ is Hilbert--Schmidt self-adjoint so its
eigenvalues are real. Corrected and proved by the prover lane (D4).
\end{proof}

\begin{observation}[No odd-sector Alon--Boppana bound]
\label{obs:odd-gap-no-alon-boppana}
\claimstatus{obs:odd-gap-no-alon-boppana}{sketched}
The drafted odd-sector analogue of Hastings' lower bound, $\specrad(\chan|_1)\ge\lamH-o(1)$, is
false (prover, D5, refuted): the degree-four channel $\tfrac14(2\Ad1+2\Ad P)$ on
$\mathbb C^{m|m}$ has $\chan|_1 = 0$ for every $m$, and there are primitive degree-sixteen
families with $\specrad(\chan|_1) = 0$ and unbounded $m$. Hastings' argument
(\texttt{prov:hastings07-alon-boppana}, \texttt{0706.0556:sd10.tex:284-287},
{\small\texttt{\detokenize{N(0,m) \geq c \times (2\sqrt{D-1})^m/(m+1)^{3/2},}}})
needs every word trace $|\Tr U_w|^2\ge0$; on the odd sector word traces can be negative. The
even letter $P$ is exactly the parity jump of \cref{prop:parity-jump-uniform-shift}, which shifts
the whole odd spectrum: the odd gap has no universal lower bound, and ``Ramanujan'' for the odd
sector is an upper bound (the band), not an optimality statement.
\end{observation}

\subsection{Graded Harrow expanders}

\begin{theorem}[Index-two gradings of Harrow channels]
\label{thm:graded-harrow}
\claimstatus{thm:graded-harrow}{sketched}
Let $\chan$ be the graded representation channel of \cref{def:index-two-grading}. (i) $\Gdb =
\Ad(P)$ is a $G$-intertwiner of $\pi\otimes\bar\pi = \operatorname{End}(V)$; as $G_0$-modules the
even sector is $\pi_+\otimes\bar\pi_+\oplus\pi_-\otimes\bar\pi_-$ and the odd sector
$\pi_+\otimes\bar\pi_-\oplus\pi_-\otimes\bar\pi_+$. (ii) The trivial representation occurs once, in
the even sector (the scalars, the fixed point), and the sign representation $\epsilon$ once, in the
even sector, spanned by $P$, with $\chan(P) = \epsilon(W_S)P$. (iii) Every other constituent $\constit$
of either sector satisfies $|\constit(W_S)|\le\lamtwo(\mathrm{Cay}(G,S))$ (\cref{cit:harrow-transfer});
if $\mathrm{Cay}(G,S)$ is Ramanujan, $\chan$ is a graded Ramanujan quantum expander and, by
\cref{cor:graded-band-circle}, the odd sector of its graded Hashimoto operator carries zeros on
$|\uz| = \qs^{-1/2}$, $\qs = |S|-1$. (iv) Character formulas: $\Tr\chan^n = |S|^{-n}\sum_{|w| = n}
|\chi_\pi(w)|^2$ and $\str\chan^n = |S|^{-n}\sum_{|w| = n}|\chi_+(w)-\chi_-(w)|^2$, where $\chi_\pi =
\chi_++\chi_-$ on $G_0$ and $\chi_\pi = 0$ off $G_0$; only words with product in $G_0$ contribute to
either ring norm. (v) Artin factorisation: with $a_{k,\constit}$ the multiplicity of $\constit$ in sector $k$,
the graded quantum Ihara zeta is $\prod_\constit\det(1-\uz\Hash_\constit)^{a_{1,\constit}-a_{0,\constit}}$,
$\Hash_\constit$ the Hashimoto operator built from the individual letters $\constit(s)$ and the
reversal (not from the averaged matrix $\constit(W_S)$ alone), so the $L$-functions of the graded
expander are the $\constit$-factors, with net sign by sector.
\end{theorem}
\begin{proof}
(i) $\Ad(P)\Ad(\pi(g)) = \Ad(\epsilon(g)\pi(g)P) = \Ad(\pi(g))\Ad(P)$. (ii) $\pi$ irreducible gives
$\operatorname{End}_G(V) = \mathbb C$; $P$ spans the $\epsilon$-isotypic part since $\Ad\pi(g)P =
\epsilon(g)P$, and $\epsilon\ne1$ so it is not the trivial constituent; both are even. (iii) is
Harrow's argument on each sector. (iv) $\Tr\Ad(U) = |\Tr U|^2$ and $\str\Ad(U) = |\Tr(PU)|^2$;
$\Tr(P\pi(g)) = \chi_+(g)-\chi_-(g)$ for $g\in G_0$ and $0$ for block-off-diagonal $\pi(g)$; the
sums run over labelled words with product in $G_0$. (v) decompose each sector into $G$-isotypic
components; $\Hash$ commutes with the decomposition. Corrected and proved by the prover lane (D6,
ledger L06--L09, L37).
\end{proof}

\begin{numerical}[Three graded Ramanujan expanders]
\label{num:graded-ramanujan}
\claimstatus{num:graded-ramanujan}{numerical}
\texttt{scripts/graded\_ramanujan.py} (\texttt{outputs/graded\_ramanujan.txt}, $81$ checks).
\textbf{(a) Pauli.} $V = \mathbb C^{1|1}$, $P = Z$, letters $X,X,Y,Y,Z,Z$ (each self-adjoint letter
twice so the reversal is fixed-point-free), $\Dk = 6$, $\qs = 5$. Even spectrum $\{6,-2\}$, odd
$\{-2,-2\}$, all nontrivial in the band $|\lambda|\le2\sqrt5$ (the direct, non-lifted transfer has
$Z_P = (1+2\uz)/(1-6\uz)$ with its zero at $-2$, off the circle: the curve appears only after the
non-backtracking lift); graded Ihara--Bass on both sectors;
odd edge eigenvalues on $|\edgeev| = \sqrt5$; $\str\Hash^n = 8,32,104,640,3208,15392 = 1+5^n-
(\edgeev^n+\bar\edgeev^n)$ with $\edgeev = -1+2i$, verified against the word sum
$\sum_w|\Tr(PU_w)|^2$ for $n\le3$; graded quantum Ihara zeta $(1+2\uz+5\uz^2)/((1-\uz)(1-5\uz))$.
\textbf{(b) The 06h elliptic tensor} ($\qs = 5$, $\pmode = 1$): $\sum_s\epsilon_s\Kr{s}\Kr{s}^\dagger
= 1$, $\Edb(P) = P$, even $\{5,1\} = \trivS$, odd $\{\pi,\bar\pi\}$ on the circle with
$\Esec{1}^\dagger\Esec{1} = 5$, graded Weil--Hodge inequality for $n\le8$.
\textbf{(c) Graded Weil--LPS.} $G = \mathrm{PGL}_2(\mathbb F_p)$, $\pi$ the principal series
$I(\chi,\chi^{-1})$ with $\chi$ of order $4$ (dimension $p+1$, splitting on $\mathrm{PSL}_2$ into
halves of dimension $(p+1)/2$, $P$ found as the traceless element of the commutant of $\pi(G_0)$ by
exact group averaging), $S$ the $\qs+1$ LPS quaternions of norm $\qs$ with $(\qs/p) = -1$, all in
the odd coset; $(p;\qs) = (5;13)$ and $(13;5)$. Cayley graphs on $120$ and $2184$ vertices,
bipartite, Ramanujan (max nontrivial $|\lambda| = 4.00$ and $4.25$ against $2\sqrt\qs = 7.21$ and
$4.47$); both sectors of $\chan$ in the band, $P$-mode $-(\qs+1)$, Harrow containment of both sector
spectra in the Cayley spectrum; $\str\chan^n\ge0$ with the character formula on words of length
$\le3$, vanishing for odd $n$; graded Ihara--Bass on both sectors of the $504$- and $1176$-dimensional
Hashimoto operators; odd edge eigenvalues $36$ and $196$ on $|\edgeev| = \sqrt\qs$ (raw counts) and the rest at
$|\edgeev| = 1$; even ones on the circle plus $\{\pm\qs,\pm1\}$. The sectors share eigenvalues,
and only the net divisor survives: for $(5;13)$ there are $16$ cancelled pairs and the reduced
zeta is $(1-4\uz+13\uz^2)(1+4\uz+13\uz^2)/[(1-\uz)(1-13\uz)(1+\uz)(1+13\uz)]$ (numerator degree
$4$, no nontrivial poles); for $(13;5)$ the reduced numerator and denominator both have degree
$144$, the denominator carrying $140$ nontrivial pole factors on the circle (e.g.\ $\lambda = 0$
has net even multiplicity $4$). My first printout claimed numerator degrees $36$ and $196$ over four
trivial factors; the prover's independent audit (\texttt{finite/harrow\_audit.py}, $16$ checks)
corrected it and the script now prints the net divisor. In the $(13;5)$ case the extremal Hecke
eigenvalue $4.2497$ lives in the odd sector.
\textbf{(d) Continuum.} The jump Lindbladian of (a) commutes with $\Gdb$; even $\{0,-2(g_X+g_Y)\}$,
odd $\{-2(g_Z+g_Y),-2(g_Z+g_X)\}$; $\str(1+\epsilon\Tcm)^{L/\epsilon}\to\str e^{L\Tcm}$ at
$\epsilon = 10^{-3}$; $\str e^{t\Tcm}\ge0$.
\end{numerical}

\begin{proposition}[The qubit with unitary letters]
\label{prop:qubit-graded-zeta}
\claimstatus{prop:qubit-graded-zeta}{sketched}
For $\Dk$ unitary homogeneous letters closed under adjoint on $\mathbb C^{1|1}$, $\qs = \Dk-1$, write
the even letters $\operatorname{diag}(a_i,d_i)$ and the odd ones with off-diagonal entries $a_i,d_i$
(all of modulus one), and put $\pmode = N_{\mathrm{even}}-N_{\mathrm{odd}}$, $x =
\sum_{\mathrm{even}}a_i\bar d_i\in\mathbb R$, $y = \sum_{\mathrm{odd}}a_i\bar d_i$. Then
$\operatorname{spec}\Sig_0 = \{\Dk,\pmode\}$, $\Sig_1 = \bigl(\begin{smallmatrix}x&y\\
\bar y&x\end{smallmatrix}\bigr)$ with eigenvalues $\lambda_\pm = x\pm|y|$, and
\[
  \frac{(1-\lambda_+\uz+\qs\uz^2)(1-\lambda_-\uz+\qs\uz^2)}{(1-\uz)(1-\qs\uz)(1-\pmode\uz+\qs\uz^2)}
\]
is the graded quantum Ihara zeta. The sector Ramanujan condition is $|\pmode|,|\lambda_\pm|\le
2\sqrt\qs$ (or, all letters odd, $\pmode = -\Dk$, $x = 0$ and $|y|\le2\sqrt\qs$ after
removing the even $\pm\Dk$ modes); the divisor condition tests only uncancelled quadratics. For the
Pauli letters $\pmode = \lambda_\pm = -2$, one odd quadratic cancels the $P$-mode pole pair, and the
result is the zeta function of the elliptic curves $y^2 = x^3+4x+b$, $b\in\{0,1,4\}$, over
$\mathbb F_5$ ($8$ points, $a_5 = -2$): the periodic non-backtracking ring norms of a qubit with
Pauli letters count the points of an elliptic curve. The numerator is in general not integral (a
Ramanujan example with $N_P(1) = 16/3$), need not be a Weil polynomial, and an integral Weil
numerator need not be a curve numerator ($(1-3\uz)^4$ over $\mathbb F_9$ would give $-2$ points);
ring norms of unitary letters are weighted squared amplitudes, not word counts, unless the letters
are $0/1$ or the words are counted with multiplicity as here (prover, D13, ledger L25--L27).
\end{proposition}
\begin{proof}
\cref{thm:graded-qihara-bass} with $n_0 = n_1 = 2$; the even sector is spanned by $1$ and $P$
(\cref{prop:p-mode}); the odd block is computed in the basis $(E_{01},E_{10})$. Corrected and
proved by the prover lane (D13). The Pauli case is \cref{num:graded-ramanujan}(a); it is an
ordinary representation of the Pauli group with $G_0 = \langle iI,Z\rangle$, projective on the
abelian quotient (ledger L09).
\end{proof}
